% Missing operative details must be checked against the study records.
Animal care approval and animal characteristics were \verifytext{ breed, sex, age, and weight]}. Anaesthesia and perioperative care were \verifytext{[recorded regimen, airway management, monitoring, analgesia, and experimental endpoint]}.

A simulated cleft was created by \verifytext{scalpel incision, etc. medical terminology}, with oral exposure provided by \verifytext{retractor}. Repair comprised \verifytext{dissection, muscle mobilization, and closure layers}, using \verifytext{suture material, size, needle, and technique}. The robotic platform was the da Vinci Research Kit (dVRK)~\cite{kazanzides2014dvrk}. \Cref{fig:workflow} outlines the procedure and annotation stages.
A single surgeon performed all stages of each procedure. The analysis included seven procedures: three with 3 mm instruments and four with 8 mm instruments. 
% \verifytext{Dissection used round-tip scissors with 8 mm instruments and a custom robotic elevator with 3 mm instruments.} 
